\section{引言}
\label{sec:intro}

格点算符的重正化是从数值模拟获得诸多物理结果（如介子衰变常数、形状因子、结构函数、
混合振幅等）的必要环节。本文研究格点 $QCD$ 中复合算符的重正化。我们提出一种完全避免
使用格点微扰论的方法，允许对任意复合算符的重正化常数作非微扰确定。作为说明，我们考虑
对两费米子算符矩阵元的一些具体应用。该途径在那些无法用手征 Ward 恒等式非微扰地确定重正化
常数的情形下特别有用，例如标量密度或赝标密度\footnote{在下文中，我们把利用 Ward 恒等式确定
重正化常数的方法称为 Ward 恒等式方法。}\cite{bo}--\cite{wi}。此外，我们的建议还可应用于许多
其他情形，如有效弱哈密顿量中四费米子算符的重正化，以及重夸克有效理论中的重轻夸克流。
初步结果已在文献 \cite{dallas} 中报告。类似尝试（仅限于 $\Delta I=1/2$ 企鹅算符的发散部分）
可见文献 \cite{ber1}。

重正化的格点算符必须在无限格点截断的极限下对应于有限的手征协变算符，并满足与连续统中相同
的重正化条件。除少数例外，人们通常用格点微扰论计算格点算符的重正化常数。然而，与低维算符
混合的问题要求对所有幂次发散作非微扰减除 \cite{bo},\cite{te}--\cite{martsukuba}。
撇开这一特殊情形不谈，只要格点间距 $a$ 足够小，即 $ a^{-1} \gg \Lambda_{QCD}$，
在计算乘性常数与无量纲混合系数时使用微扰论是合理正当的。然而，在那些可以用非微扰方法检验
微扰论结果（以裸格点耦合常数 $\asl=g_{0}^{2}/4 \pi$ 作为展开参数）的情形，一般存在显著偏差。
针对这一问题迄今已提出若干解决办法：
\begin{itemize}
\item 为改进微扰级数的收敛性，文献 \cite{lepage} 的作者提出了``蝌蚪改进''（tadpole improved）
微扰论。尽管如此，对高阶贡献的无知仍然是提取物理结果时不确定性的重要来源。
\item 在某些有限情形——对应于有限算符，即矢量流与轴矢流，以及赝标密度与标量密度重正化常数
之比——可以借助手征 Ward 恒等式 \cite{bo}--\cite{wi} 完全非微扰地确定重正化常数。
\item 可以直接在强子矩阵元上施加非微扰重正化条件。文献 \cite{gav1} 用这一办法减除
$\Delta I=1/2$ 算符的发散部分。代价是：每在强子态上施加一个减除条件，就要牺牲一个物理预言。
因此当需要若干个重正化条件时，会损失大量预言能力 \cite{bo,te}。
\end{itemize}

{\itshape 我们的建议是：在固定规范中、给定具有大虚度的离壳外态的条件下，直接在夸克与胶子格林
函数上以非微扰方式施加重正化条件。}

方法的实质是模仿微扰论中通常的做法：通过令某些适当选取的、在固定规范中于离壳夸克和胶子外态
之间计算的格林函数等于其树图值，来固定某算符的重正化条件。例如，对于一般的两夸克算符
$O_\Gamma= \bar \psi \Gamma \psi$，可施加条件
\beq Z_\Gamma \langle p \vert O_\Gamma \vert p \rangle\vert_{p^2=-\mu^2}=
\langle p \vert O_\Gamma \vert p \rangle_0, \eeq
其中 $\Gamma$ 是某个狄拉克矩阵，$\langle p \vert O_\Gamma \vert p \rangle_0$
是树图级矩阵元。向更复杂情形（包括四费米子算符与算符混合）的推广是直接的。只要都以同一个
重正化耦合常数表达，这一流程就在所有正规化方案中定义了同样的重正化算符，即同样的 Wilson
系数函数。不过此时系数函数依赖于外部态与规范，二者都必须明确指明。

原则上，这一方案完全解决了格点微扰论中大修正的问题，因为这些修正被自动纳入重正化常数之中。
将格点算符与在连续统重正化方案（为明确起见取 $\overline{MS}$ 方案）中重正化的相应算符相匹配，
便只需连续统微扰论。当次领头修正已知时，物理矩阵元的误差为连续统微扰展开中的 $O((\as)^2)$，
而连续统展开预期收敛性更好。匹配条件的连续统微扰计算是所有方案（标准或``蝌蚪改进''微扰论、
非微扰重正化）的共同步骤。与只适用于少数（尽管十分重要）情形的 Ward 恒等式方法不同，我们的
建议可应用于任意复合算符\footnote{仍然正确的是：在 Ward 恒等式方法可用的场合，其统计精度
一般略好一些，见第 \ref{sec:nonper} 节。}。

在固定规范中使用夸克态之间格林函数的可行性，已在文献 \cite{wi} 中于使用 Ward 恒等式作非微扰
重正化的语境下讨论过。事实上就轴矢流而言，文中表明信号相当好，且结果与相应的规范不变的强子
测定一致。

至于重正化条件，只要能够固定外部态的虚度 $\mu$ 并满足条件
$ \Lambda_{QCD} \ll \mu \ll 1/a $，以便同时控制非微扰效应与离散化效应，我们的建议就预期有效。
我们强调这一要求为所有方法所共有。我们的数值研究结果令人鼓舞：它们表明，在当前格点模拟所用
格点耦合常数的数值下，确实存在一个可以应用本方法的 $\mu$ ``窗口''。

本文其余部分的安排如下：下一节以对数发散的两夸克算符（如标量密度与赝标密度）为例，给出定义
非微扰方法所需的基本公式；第 3 节把讨论推广到所有复合算符；随后一节讨论两夸克算符方法在数值
模拟中的实现；第 5 节给出有限体积格点上重正化常数微扰计算的细节；第 6 节给出数值模拟得到的
重正化常数的数值，并与微扰论结果比较；最后一节给出结论、数值结果小结以及对本文思想可能进一
步应用的讨论。
